% ─────────────────────────────────────────────────
%  AVANT-PROPOS
% ─────────────────────────────────────────────────
\chapter*{Avant-propos}
\phantomsection
\addcontentsline{toc}{chapter}{Avant-propos}
\markboth{Avant-propos}{Avant-propos}

L'École Supérieure Polytechnique de Dakar (\textbf{ESP}), rattachée à
l'Université Cheikh Anta Diop (\textbf{UCAD}), assure des formations
professionnelles et d'ingénieurs dans plusieurs domaines scientifiques,
techniques et de gestion \cite{esp_officiel}. Le présent mémoire s'inscrit dans
la formation suivie au Département Génie Informatique, qui associe les acquis
académiques à une expérience pratique en entreprise.

Dans ce cadre, le projet de fin d'études permet de conduire une étude appliquée,
depuis l'analyse des besoins jusqu'à la réalisation et à l'évaluation d'une
solution informatique dans son contexte professionnel.

Ainsi, nous avons effectué un stage au sein de la structure
\textbf{INNOV4AFRICA} et travaillé sur le sujet suivant~:
\textit{Conception d'un système d'information bancaire intelligent~: cas des
modules Relation Client et SGI (Société de Gestion et d'Intermédiation)}.

\cleardoublepage
